\documentclass[12pt]{article}

% --------------------
% PAQUETES BÁSICOS
% --------------------
\usepackage[spanish,es-noshorthands]{babel}
\usepackage[utf8]{inputenc}      % Si usas pdflatex
\usepackage[T1]{fontenc}
\usepackage{lmodern}

\usepackage{amsmath, amssymb, amsthm}
\usepackage{geometry}
\usepackage{booktabs}   
\usepackage{setspace}
\usepackage{graphicx}
\usepackage{tikz}
\usepackage{tikz-cd}
\usepackage[
  colorlinks=true,
  linkcolor=blue,
  citecolor=blue,
  urlcolor=red
]{hyperref}

\newtheorem{thm}{Teorema}[section]
\newtheorem{dfn}[thm]{Definición}
\newtheorem{lem}[thm]{Lema}
\newtheorem{ej}[thm]{Ejemplo}
\newtheorem{cor}[thm]{Corolario}
\newtheorem{obs}[thm]{Observación}
\newtheorem{con}[thm]{Conjetura}
\newtheorem{prop}[thm]{Proposición}

\DeclareMathOperator{\Frm}{Frm}
\DeclareMathOperator{\Fil}{Fil}
\DeclareMathOperator{\Ord}{Ord}
\DeclareMathOperator{\Top}{Top}
\DeclareMathOperator{\Sp}{sp}
\DeclareMathOperator{\pt}{pt}
\DeclareMathOperator{\id}{id}
\DeclareMathOperator{\tp}{tp}
\DeclareMathOperator{\fd}{fd}
% --------------------
% CONFIGURACIÓN
% --------------------
\geometry{margin=2.5cm}
\onehalfspacing

% --------------------
% DATOS DEL DOCUMENTO
% --------------------
\title{Reporte mensual de actividades}
\date{1-abril-2026 a 30-abril-2026}
\author{Juan Carlos Monter Cortés}

% --------------------
% INICIO DEL DOCUMENTO
% --------------------
\begin{document}

\maketitle
% \tableofcontents

\section{Actividades realizadas}
Uno de nuestros objetivos es demostrar que la siguiente afirmación es cierta.
\[
\text{\textbf{Afirmación:} Si } A\in \Frm \text{ es } T_1 \text{ y apilado, entonces } A \text{ es espacial.}
\]
Para realizarlo debemos verificar que el morfismo reflexión espacial es inyectivo, es decir, para $a, b\in A$ con $a\neq b$, se cumple que
\[
U(a)\nsubseteq U(b)\; \iff \; \exists\, p\in \pt(A) \text{ tal que } a\leq p \text{ y } b\nleq p.
\]

Notemos que para $a\nleq b$ como antes, se cumple que $\uparrow a$ es un filtro en $A$ y $\downarrow b$ es un ideal. Además, se cumple que 
\[
\uparrow a \cap \downarrow b = \emptyset.
\]

Denotamos por $\Fil(A)$ al conjunto de todos los filtros en $A$. Si consideremos el conjunto
\[
\mathcal{F}(a,b)=\{F\in \Fil(A)\mid a\in F\text{ y }b\notin F\},
\]
entonces por Lema de Zorn se cumple que existe un filtro maximal $\mathcal{M}\in \mathcal{F}(a,b)$. Ya que $\mathcal{M}\in \Fil(A)$, podemos construir el núcleo ajustado que admite a $\mathcal{M}$, esto es
\[
\mathcal{M}\subseteq \nabla(v_{\mathcal{M}})
\]
donde $v_\mathcal{M}$ es el núcleo ajustado. Dado que $v_\mathcal{M}\in NA$, tenemos el marco cociente $A_{v_{\mathcal{M}}}$ cuyo espacio de puntos está dado por
\[
\pt(A_{v_\mathcal{M}})=S\setminus \mathcal{M}
\]
para $S=\pt A$. Notemos que $S\setminus \mathcal{M}$ no es vacío, pues en caso contrario
\[
S\setminus \mathcal{M}=\emptyset\Rightarrow M\subseteq S
\]
lo cual es una contradicción ya que $\mathcal{M}$ es un filtro y $1\in \mathcal{M}$ y $1\notin S$. Como consecuencia, existe  $p\in S$ tal que $p\notin \mathcal{M}$.
Luego, si $p\notin \mathcal{M}$, entonces $a\nleq p$, pues si $a\leq p$ eso implicaría que $p\in \mathcal{M}$. Lo cual no puede ocurrir.

\bibliographystyle{amsalpha}
\bibliography{research2}


\end{document}